\chapter{Testy parametryczne: średnie, wariancje, częstości i ANOVA}
\label{ch:parametric-tests}
\paragraph{Powiązania i odsyłacze.} Testy opierają się na rozkładach statystyk z rozdziału~\ref{ch:statystyki} i konstrukcji testów z rozdziału~\ref{ch:testing}.


\section{Jedna średnia: kiedy test normalny, a kiedy t-Studenta?}
Rozpocznijmy od pomiarów czasu pracy urządzenia. Interesuje nas,
czy średni czas przekracza wartość referencyjną $\mu_0$.
Jeśli znamy odchylenie standardowe modelu normalnego, statystyka
$Z=\sqrt n(\bar X-\mu_0)/\sigma$ ma dokładnie rozkład
$\Normal(0,1)$ pod hipotezą $H_0:\mu=\mu_0$.
Najczęściej $\sigma$ nie jest jednak znane i trzeba je oszacować
z tych samych danych. Wtedy pojawia się rozkład t-Studenta.

\begin{twierdzenie}[Dokładny test t dla jednej średniej]
\label{thm:one-sample-t}
Niech $n\ge2$ i $X_1,\ldots,X_n\iid\Normal(\mu,\sigma^2)$,
gdzie $\mu\in\R$, $\sigma>0$. Niech
$\bar X=n^{-1}\sum_iX_i$ i
$S^2=(n-1)^{-1}\sum_i(X_i-\bar X)^2$.
Pod $H_0:\mu=\mu_0$ statystyka
\[
 T=\frac{\bar X-\mu_0}{S/\sqrt n}
\]
ma rozkład t-Studenta z $n-1$ stopniami swobody, z wyjątkiem
zdarzenia $S=0$, które ma prawdopodobieństwo zero.
Tutaj $t_{\nu,q}$ oznacza kwantyl rzędu $q\in(0,1)$
rozkładu t-Studenta z $\nu>0$ stopniami swobody.
Dla $0<\alpha<1/2$ test dwustronny odrzucający dla
$|T|>t_{n-1,1-\alpha/2}$ ma poziom dokładnie $\alpha$,
a jednostronny dla $T>t_{n-1,1-\alpha}$ również ma
poziom dokładnie $\alpha$ przy $H_0:\mu\le\mu_0$.
\end{twierdzenie}
\begin{proof}
Z rozkładu próby normalnej wyprowadzonego w
rozdziale~\ref{ch:statystyki} mamy
$Z=\sqrt n(\bar X-\mu)/\sigma\sim\Normal(0,1)$ oraz
$V=(n-1)S^2/\sigma^2\sim\chi^2_{n-1}$, przy czym
$Z$ i $V$ są niezależne. Pod $H_0$ statystyka $T$ jest
równa $Z/\sqrt{V/(n-1)}$, co definiuje rozkład t-Studenta.
Dla alternatywy jednostronnej przy $\mu\le\mu_0$
rozkład $T$ jest przesunięty przez niedodatni składnik
$\sqrt n(\mu-\mu_0)/\sigma$ w liczniku, przy zachowanej
niezależności mianownika. Stąd prawdopodobieństwo przekroczenia
ustalonego dodatniego progu jest największe przy $\mu=\mu_0$.
Symetria i ciągłość rozkładu t zapewniają dokładny poziom
dla obu reguł.
\end{proof}

\begin{uwaga}[Normalność czego?]
Test t dla jednej średniej ma dokładny rozkład pod
założeniem, że \emph{obserwacje} są iid normalne. W teście t dla
różnic sparowanych normalności wymagamy od \emph{różnic}, nie
koniecznie od każdej serii oddzielnie. W regresji i ANOVA
normalność dotyczy rozkładu błędów warunkowych, a nie połączonej
próby ze wszystkich grup. Przy dużych próbach przybliżenia
na podstawie CLT mogą działać bez normalności danych, ale
zależą od ogonów rozkładu, skończoności wariancji i schematu
próby.
\end{uwaga}

\begin{stwierdzenie}[Przedział dokładny dla średniej normalnej]
Pod założeniami twierdzenia~\ref{thm:one-sample-t}
przedział losowy
\[
 \left[\bar X-t_{n-1,1-\alpha/2}\frac S{\sqrt n},\;
 \bar X+t_{n-1,1-\alpha/2}\frac S{\sqrt n}\right]
\]
obejmuje $\mu$ z prawdopodobieństwem dokładnie $1-\alpha$.
\end{stwierdzenie}
\begin{proof}
To zdarzenie jest równoważne
$|\sqrt n(\bar X-\mu)/S|\le t_{n-1,1-\alpha/2}$.
Twierdzenie~\ref{thm:one-sample-t} dla prawdziwej średniej
zapewnia prawdopodobieństwo $1-\alpha$.
\end{proof}

\section{Pomiary powtarzane: różnice par zamiast dwóch prób niezależnych}
\begin{twierdzenie}[Test t dla par]\label{thm:paired-t}
Niech pary $(X_i,Y_i)$, $i=1,\ldots,n$, będą niezależne
między jednostkami, a różnice $D_i=X_i-Y_i$ mają wspólny
rozkład $\Normal(\delta,\tau^2)$ z $\tau>0$ i $n\ge2$.
Wtedy przy $H_0:\delta=0$ statystyka
$T_D=\sqrt n\bar D/S_D$ ma dokładnie rozkład $t_{n-1}$.
Nie wymagamy niezależności $X_i$ i $Y_i$ wewnątrz pary.
\end{twierdzenie}
\begin{proof}
Zmienne $D_i$ tworzą próbę prostą normalną, więc jest to
bezpośrednie zastosowanie twierdzenia~\ref{thm:one-sample-t}
do obserwacji $D_1,\ldots,D_n$.
\end{proof}

\begin{zadanie}[Przed i po zmianie procedury]
Dla ośmiu osób wyniki przed zmianą wynoszą
$54,58,57,60,56,55,61,59$, a po zmianie
$51,56,53,57,55,53,56,57$. Czy średnia różnica jest dodatnia?
\par\emph{Rozwiązanie.} Różnice to
$3,2,4,3,1,2,5,2$, stąd
$\bar d=2{,}75$, $s_D^2=23/14\approx1{,}6429$,
$T_D=\sqrt8\,\bar d/s_D\approx6{,}068$,
przy siedmiu stopniach swobody. W teście dwustronnym
wartość p jest około $0{,}00051$; w zaplanowanym teście
jednostronnym na spadek wyniku po interwencji jest około
$0{,}00025$. To jednak nie dowodzi przyczynowości, jeżeli
dane nie pochodzą z odpowiednio zaprojektowanego eksperymentu.
\end{zadanie}

\section{Dwie niezależne grupy: Student czy Welch?}
\begin{twierdzenie}[Dwie średnie przy wspólnej wariancji]
\label{thm:two-sample-t}
Niech $X_1,\ldots,X_n\iid\Normal(\mu_X,\sigma^2)$ i
$Y_1,\ldots,Y_m\iid\Normal(\mu_Y,\sigma^2)$ będą niezależnymi
próbami, gdzie $n,m\ge2$, a $\sigma>0$ jest nieznane.
Niech
$S_p^2=[(n-1)S_X^2+(m-1)S_Y^2]/(n+m-2)$.
Pod $H_0:\mu_X-\mu_Y=\delta_0$ zachodzi
\[
 T_p=\frac{(\bar X-\bar Y)-\delta_0}
 {S_p\sqrt{1/n+1/m}}\sim t_{n+m-2}.
\]
\end{twierdzenie}
\begin{proof}
Ponieważ próby są niezależne,
$Z=[(\bar X-\bar Y)-(\mu_X-\mu_Y)]/
[\sigma\sqrt{1/n+1/m}]$ ma rozkład $\Normal(0,1)$.
Dla każdej próby z osobna resztowa suma kwadratów podzielona
przez $\sigma^2$ jest zmienną $\chi^2$ o odpowiednio
$n-1$ i $m-1$ stopniach swobody; obie są niezależne
między sobą i od obu średnich. Ich suma
$(n+m-2)S_p^2/\sigma^2$ ma rozkład
$\chi^2_{n+m-2}$. Iloraz $Z$ i pierwiastka ze
średniej tej sumy daje dokładny rozkład t.
\end{proof}

Gdy wariancje są różne, podstawienie wariancji wspólnej
może istotnie zniekształcić test, szczególnie jeśli
liczebności grup też się różnią. Wtedy naturalną alternatywą
jest \emph{test Welcha}:
\[
 T_W=\frac{(\bar X-\bar Y)-\delta_0}
 {\sqrt{S_X^2/n+S_Y^2/m}},\qquad
 \nu=\frac{(S_X^2/n+S_Y^2/m)^2}
 {(S_X^2/n)^2/(n-1)+(S_Y^2/m)^2/(m-1)}.
\]
Rozkład $T_W$ jest zwykle \emph{przybliżany} rozkładem
$t_\nu$, gdzie $\nu$ to losowe, niekoniecznie całkowite
przybliżenie liczby stopni swobody Satterthwaite'a.
Nie jest to twierdzenie o dokładnym rozkładzie t przy
nierównych wariancjach, nawet jeśli obie próby są normalne.

\begin{zadanie}[Dwie technologie produkcji]\label{zad:welch-case}
Czasy obróbki w technologii A to
$18,19,20,19,21,20,18,22$; w technologii B:
$22,23,21,24,22,25,23,24$. Porównaj średnie.
\par\emph{Rozwiązanie.} Mamy $n=m=8$,
$\bar x=19{,}625$, $\bar y=23$, $s_X^2\approx1{,}9821$,
$s_Y^2\approx1{,}7143$. Statystyka Welcha wynosi
$T_W\approx-4{,}9651$, $\nu\approx13{,}927$,
a dwustronna wartość p około $0{,}00021$.
W tym przykładzie test ze wspólną wariancją daje bardzo
zbliżony wynik, ale to \emph{nie dowodzi}, że w dowolnym
niezrównoważonym badaniu obie procedury są równoważne.
Do raportu należy dołączyć różnicę średnich
$\bar x-\bar y=-3{,}375$ i jej przedział ufności.
\end{zadanie}

\begin{uwaga}[Nie wybieramy testu wyłącznie przez pretest wariancji]
Procedura „najpierw test równości wariancji, potem Student lub
Welch zależnie od p-value” zmienia rozkład końcowej reguły i
nie jest automatycznie testem o nominalnym poziomie.
Gdy nie mamy silnego powodu, by zakładać równość wariancji,
Welch jest zwykle rozsądnym punktem wyjścia dla dwóch
niezależnych grup. Z kolei silna asymetria, ciężkie ogony,
zależność obserwacji czy mała próba wymagają dodatkowej oceny:
sam Welch nie usuwa wszystkich tych problemów.
\end{uwaga}

\section{Testy wariancji i kiedy nie należy ich używać}
\begin{twierdzenie}[Test chi-kwadrat jednej wariancji]
Niech $X_i\iid\Normal(\mu,\sigma^2)$, $n\ge2$,
$\mu\in\R$ nieznana i $\sigma>0$. Niech $v_0>0$. Pod $H_0:\sigma^2=v_0$
statystyka $Q=(n-1)S^2/v_0$ ma rozkład
$\chi^2_{n-1}$. Dwustronny test odrzuca dla
$Q<\chi^2_{n-1,\alpha/2}$ lub
$Q>\chi^2_{n-1,1-\alpha/2}$, gdzie dolne i górne
kwantyle odpowiadają rozkładowi $\chi^2$.
\end{twierdzenie}
\begin{proof}
Wynik o rozkładzie sumy kwadratów scentrowanej próby
normalnej został udowodniony w rozdziale~\ref{ch:statystyki}.
Pod hipotezą zerową podstawiamy $\sigma^2=v_0$.
\end{proof}

\begin{twierdzenie}[Test F równości dwóch wariancji]
Niech niezależne próby normalne mają odpowiednio
liczebności $n,m\ge2$ i wariancje $\sigma_X^2,\sigma_Y^2>0$.
Pod $H_0:\sigma_X^2=\sigma_Y^2$ iloraz
$F=S_X^2/S_Y^2$ ma rozkład Fishera--Snedecora
$F_{n-1,m-1}$.
\end{twierdzenie}
\begin{proof}
Z niezależności prób i rozkładu wariancji normalnej
$U=(n-1)S_X^2/\sigma_X^2\sim\chi^2_{n-1}$ i
$V=(m-1)S_Y^2/\sigma_Y^2\sim\chi^2_{m-1}$ są
niezależne. Pod $H_0$ zachodzi
$F=[U/(n-1)]/[V/(m-1)]$, co z definicji daje
rozkład Fishera--Snedecora.
\end{proof}

Rozkład dokładny jest tutaj konsekwencją normalności.
W praktyce test F jest bardzo wrażliwy na jej naruszenie.
Jeśli zadaniem jest sprawdzenie założenia jednorodności
rozproszenia w wielu grupach, często używa się procedur
Levene'a, Browna--Forsythe'a lub Flignera--Killeena.
Ich konstrukcje, ograniczenia i porównania z testami
normalności (w tym Shapiro--Wilka) omówimy osobno w części
o diagnostyce modeli.

\section{Jednoczynnikowa ANOVA: skąd pochodzi statystyka F?}
\label{sec:anova-introduction}
\begin{definicja}[Model jednoczynnikowej ANOVA]
Niech $k\ge2$ i $n_j\ge2$ dla $j=1,\ldots,k$.
Obserwacje $Y_{ji}$ są niezależne, a
\[
 Y_{ji}=\mu_j+\varepsilon_{ji},\qquad
 \varepsilon_{ji}\iid\Normal(0,\sigma^2),\qquad\sigma>0.
\]
Hipoteza zerowa ma postać
$H_0:\mu_1=\cdots=\mu_k$.
Niech $N=\sum_jn_j$,
$\bar Y_j=n_j^{-1}\sum_{i=1}^{n_j}Y_{ji}$ i
$\bar Y=N^{-1}\sum_jn_j\bar Y_j$.
Definiujemy sumy kwadratów
$SS_B=\sum_jn_j(\bar Y_j-\bar Y)^2$ (między grupami)
oraz $SS_W=\sum_j\sum_i(Y_{ji}-\bar Y_j)^2$
(wewnątrz grup).
\end{definicja}

\begin{lemat}[Ortogonalny rozkład sum kwadratów]
\label{lem:anova-decomp}
Dla dowolnych liczb rzeczywistych $y_{ji}$ przy
zdefiniowanych średnich zachodzi
\[
 \sum_{j=1}^k\sum_{i=1}^{n_j}(y_{ji}-\bar y)^2
 =\sum_j\sum_i(y_{ji}-\bar y_j)^2
 +\sum_jn_j(\bar y_j-\bar y)^2.
\]
\end{lemat}
\begin{proof}
Dla obserwacji z grupy $j$ zapisujemy
$y_{ji}-\bar y=(y_{ji}-\bar y_j)+(\bar y_j-\bar y)$.
Po podniesieniu do kwadratu i zsumowaniu wyraz mieszany
znika, gdyż $\sum_i(y_{ji}-\bar y_j)=0$.
Pozostałe dwa składniki są dokładnie $SS_W$ i $SS_B$.
\end{proof}

\begin{twierdzenie}[Dokładny rozkład statystyki ANOVA]
\label{thm:anova-f}
W modelu jednoczynnikowej ANOVA pod $H_0$ zmienne
$SS_B/\sigma^2$ i $SS_W/\sigma^2$ są niezależne,
a ich rozkłady wynoszą odpowiednio $\chi^2_{k-1}$
oraz $\chi^2_{N-k}$. W konsekwencji
\[
 F=\frac{SS_B/(k-1)}{SS_W/(N-k)}\sim F_{k-1,N-k}
\]
i test odrzucający przy dużych wartościach $F$ ma
poziom dokładnie $\alpha$, jeśli użyjemy kwantyla
$F_{k-1,N-k,1-\alpha}$.
\end{twierdzenie}
\begin{proof}
Zapiszmy wszystkie $N$ obserwacji jako wektor
$Y\in\R^N$ o postaci $Y=\mu\mathbf1+\varepsilon$
pod $H_0$, gdzie $\varepsilon\sim\Normal_N(0,\sigma^2I_N)$.
Niech $G\subset\R^N$ oznacza podprzestrzeń wektorów
stałych w obrębie każdej grupy. Ma ona wymiar $k$ i zawiera
wektor $\mathbf1$. Oznaczmy przez $P_G$ rzut
ortogonalny na $G$, a przez $P_0$ rzut na
$\operatorname{span}\{\mathbf1\}$. Wówczas
$P_G-P_0$ i $I_N-P_G$ są rzutami ortogonalnymi na
wzajemnie prostopadłe podprzestrzenie o wymiarach
odpowiednio $k-1$ i $N-k$.

Wektor $P_GY$ ma na pozycjach grupy $j$ wartość
$\bar Y_j$, podczas gdy $P_0Y=\bar Y\mathbf1$.
Zatem $SS_B=\|(P_G-P_0)Y\|^2$ i
$SS_W=\|(I_N-P_G)Y\|^2$.
Oba rzuty zerują $\mathbf1$, stąd pod $H_0$ możemy w tych
wyrażeniach zastąpić $Y$ przez $\varepsilon$.
Po wyborze ortonormalnych baz obu podprzestrzeni
współrzędne rzutów $\varepsilon/\sigma$ są niezależnymi
standardowymi zmiennymi normalnymi. Wynika to z
ortogonalnej niezmienniczości $\Normal_N(0,I_N)$:
współrzędne obu rzutów są funkcjami rozłącznych zbiorów
niezależnych współrzędnych. Sumy ich kwadratów mają
zatem wskazane rozkłady chi-kwadrat i są niezależne.
Iloraz niezależnych średnich takich sum ma z definicji
rozkład $F_{k-1,N-k}$.
\end{proof}

\begin{zadanie}[Trzy metody produkcji]\label{zad:anova-three}
Dla trzech niezależnych grup po osiem obserwacji mamy:
$A=(18,19,20,19,21,20,18,22)$,
$B=(22,23,21,24,22,25,23,24)$ i
$C=(25,27,26,24,28,26,27,25)$.
Przetestuj równość średnich i zinterpretuj wynik.
\par\emph{Rozwiązanie.} Średnie grup wynoszą odpowiednio
$19{,}625$, $23$ i $26$. Dla $N=24$ mamy $SS_B=162{,}75$, $SS_W=37{,}875$,
zatem $F\approx45{,}119$ o stopniach swobody $(2,21)$ i
wartości p około $2{,}50\cdot10^{-8}$.
Przykładowa miara wielkości efektu to
$\eta^2=SS_B/(SS_B+SS_W)\approx0{,}811$; opisuje część
zmienności empirycznej wyjaśnianą przez podział na grupy,
nie dowodzi zaś przyczynowego wpływu metody.
Hipoteza wszystkich równych średnich zostaje odrzucona przy
konwencjonalnych poziomach istotności. Samo ANOVA nie
wskazuje jednak, \emph{które} pary średnich się różnią.
Na podstawie testu zbiorczego nie wolno też bez dodatkowej
procedury deklarować istotności wszystkich trzech porównań.
\end{zadanie}

\begin{uwaga}[Wielkość efektu a estymator poprawiony]
W jednoczynnikowej ANOVA miara $\eta^2=SS_B/(SS_B+SS_W)$
jest prostym opisem danych, lecz ma dodatnią skłonność do
zawyżania udziału wariancji w małych próbach, także gdy
wszystkie średnie populacyjne są równe. Często raportuje się
również miarę $\omega^2$, zdefiniowaną (przed ewentualnym
obcięciem ujemnego wyniku do zera) przez
\[
 \widehat{\omega}^2=
 \frac{SS_B-(k-1)MS_W}{SS_B+SS_W+MS_W},\qquad
 MS_W:=\frac{SS_W}{N-k}.
\]
Sama wielkość efektu nie koryguje błędów doboru próby ani
nie dowodzi przyczynowości. W raportach należy rozróżnić
miarę empiryczną, estymator parametru populacyjnego oraz
jego niepewność.
\end{uwaga}

\begin{uwaga}[ANOVA, Welch--ANOVA, ANCOVA i MANOVA]
Jednoczynnikowa ANOVA zakłada wspólną wariancję błędów i
niezależność obserwacji. Przy różnicach wariancji można rozważyć
ANOVA Welcha. Jeśli wynik zależy również od ciągłego
kowariantu (np. wieku albo wartości wyjściowej), w określonym
modelu można stosować ANCOVA. Jeśli każda jednostka ma kilka
współzależnych wyników, pojawia się MANOVA. Dla pomiarów
powtarzanych konieczne jest uwzględnienie zależności
wewnątrz jednostki. Każda z tych konstrukcji będzie miała
osobny, szczegółowy rozdział; nie są one zamiennymi etykietami.
\end{uwaga}

\section{Test dokładny dla częstości i testy chi-kwadrat}
\begin{stwierdzenie}[Test dwumianowy jako test dokładny]
Niech $S\sim\Bin(n,p)$, $n\in\N$, $n\ge1$, oraz
$p_0\in(0,1)$. Przy hipotezie $H_0:p\le p_0$ i
alternatywie $H_1:p>p_0$ wartość p dla obserwacji $S=s$
wynosi $P_{p_0}(S\ge s)$. Przy alternatywie
$H_1:p<p_0$ wobec $H_0:p\ge p_0$ wynosi
$P_{p_0}(S\le s)$. Odrzucanie odpowiednio dla małych
wartości p daje poziom co najwyżej $\alpha$.
\end{stwierdzenie}
\begin{proof}
Wynika to z monotoniczności stochastycznej
$\Bin(n,p)$ względem $p$, udowodnionej przez
sprzężenie w twierdzeniu~\ref{thm:binomial-ump}, oraz
z twierdzenia~\ref{thm:pvalue-valid}.
\end{proof}

\begin{twierdzenie}[Statystyka Pearsona dla wielomianowego modelu prostego]
\label{thm:pearson-chi-square}
Niech $m\ge2$ i $p_1,\ldots,p_m>0$, przy
$\sum_jp_j=1$. Niech $(N_1,\ldots,N_m)$ ma rozkład
wielomianowy z parametrami $n$ i
$(p_1,\ldots,p_m)$. Wtedy przy $n\to\infty$
\[
 Q_n=\sum_{j=1}^m\frac{(N_j-np_j)^2}{np_j}
 \xrightarrow{d}\chi^2_{m-1}.
\]
\end{twierdzenie}
\begin{proof}
Każdą z $n$ niezależnych obserwacji reprezentujemy wektorem
jednostkowym $Z_i\in\R^m$ o rozkładzie
$P(Z_i=e_j)=p_j$. Mamy
$\E Z_i=p=(p_1,\ldots,p_m)^{\mathsf T}$ oraz
$\Cov(Z_i)=\diag(p)-pp^{\mathsf T}$.
Z wielowymiarowego CLT
$\sqrt n(N/n-p)\xrightarrow{d}\Normal_m(0,
\diag(p)-pp^{\mathsf T})$.
Niech $D=\diag(p)$ i $u=(\sqrt{p_1},\ldots,\sqrt{p_m})$.
Po przekształceniu liniowym $D^{-1/2}$ macierz
kowariancji granicznej ma postać
$I_m-uu^{\mathsf T}$, czyli rzut ortogonalny na
$(m-1)$-wymiarową podprzestrzeń $u^\perp$, bo
$\|u\|^2=1$. Zgodnie z twierdzeniem o ciągłym
odwzorowaniu $Q_n$ zbiega do kwadratu normy wektora
normalnego z tą macierzą kowariancji. W ortonormalnej
bazie $u^\perp$ jest to suma kwadratów $m-1$
niezależnych zmiennych $\Normal(0,1)$, czyli
$\chi^2_{m-1}$.
\end{proof}

\begin{uwaga}[Estymowane prawdopodobieństwa i małe liczebności]
Powyższe twierdzenie dotyczy \emph{ustalonych}, znanych
wartości $p_j$ pod $H_0$. Jeśli estymujemy parametry z tej
samej próby, liczba stopni swobody może się zmienić; wymaga
to osobnej analizy regularności i rzędu modelu.
Gdy oczekiwane liczebności $np_j$ są małe, asymptotyczne
przybliżenie może być niedokładne. Reguła typu
„wszystkie oczekiwane liczebności co najmniej 5” jest
praktyczną heurystyką, a nie twierdzeniem zapewniającym
określony błąd poziomu testu. Testy dokładne i procedury
symulacyjne bywają wówczas lepsze.
\end{uwaga}

\begin{uwaga}[Poziom 5\% i dwie strony]
Test jednostronny na poziomie $\alpha$ ma dualną
\emph{jednostronną} granicę ufności poziomu $1-\alpha$.
Dwustronny przedział równych ogonów poziomu $1-\alpha$
wykorzystuje zwykle kwantyle odpowiadające $\alpha/2$.
Nie wolno utożsamiać automatycznie obu procedur. W
rozdziale~\ref{ch:confidence-intervals} przeprowadzimy
formalną konstrukcję przez odwracanie testów.
\end{uwaga}

\section{Praktyczne drzewo wyboru testu}
\begin{center}\small
\begin{tabular}{>{\raggedright\arraybackslash}p{0.31\textwidth}>{\raggedright\arraybackslash}p{0.27\textwidth}>{\raggedright\arraybackslash}p{0.31\textwidth}}
\toprule
\textbf{Dane i pytanie} & \textbf{Punkt wyjścia} & \textbf{Co trzeba dodatkowo sprawdzić?}\\
\midrule
Jedna średnia, znana $\sigma$, iid normalne & Test z & Czy $\sigma$ jest naprawdę znana z modelu?\\
Jedna średnia, nieznana $\sigma$ & Test t & Rozkład, ciężkie ogony, niezależność\\
Dwa pomiary tych samych osób & Test t dla różnic & Zależność pomiarów i rozkład \emph{różnic}\\
Dwie niezależne grupy & Test Welcha & Dobór próby, asymetria, ogony, liczebności\\
Więcej niż dwie niezależne grupy & ANOVA lub ANOVA Welcha & Wariancje, niezależność, efekty i post hoc\\
Jedna wariancja w modelu normalnym & Test $\chi^2$ & Wrażliwość na brak normalności\\
Równość dwóch wariancji w modelu normalnym & Test F & Bardzo duża wrażliwość na ogony\\
Odsetek i mała liczba sukcesów & Dokładny test dwumianowy & Schemat próby i kierunek alternatywy\\
Kilka kategorii i znane proporcje & Test Pearsona & Oczekiwane liczebności; parametry estymowane\\
Wiele pomiarów tej samej osoby & Model mieszany lub metody dla powtarzanych pomiarów & Struktura korelacji, braki danych\\
\bottomrule
\end{tabular}
\end{center}

\paragraph{Raportowanie wyników.}
Poprawny raport powinien zawierać rodzaj testu, hipotezę
zerową i alternatywną, jednostkę analizy, liczebności,
oszacowanie efektu i jego przedział, statystykę wraz ze
stopniami swobody, wartość p, sprawdzone założenia i
ograniczenia. Nie należy pisać „test potwierdził normalność”
ani „brak istotności dowodzi równości”.

\begin{lstlisting}[language=Python,caption={Kompletny przykład: Welch, ANOVA i diagnostyka w Pythonie (SciPy).}]
from scipy.stats import f_oneway, levene, shapiro, ttest_ind

A = [18, 19, 20, 19, 21, 20, 18, 22]
B = [22, 23, 21, 24, 22, 25, 23, 24]
C = [25, 27, 26, 24, 28, 26, 27, 25]

welch = ttest_ind(A, B, equal_var=False)
anova = f_oneway(A, B, C)
normality_A = shapiro(A)
variances = levene(A, B, C, center="median")

print("Welch A vs B:", welch)
print("ANOVA A, B, C:", anova)
print("Shapiro-Wilk A:", normality_A)
print("Brown-Forsythe (Levene median):", variances)
\end{lstlisting}

Ten program ilustruje sposób wywołania procedur, ale jego
wyniki wymagają \emph{analizy projektu}, a nie tylko odczytania
kolumny p-value. Shapiro--Wilk dla surowej grupy A nie jest
uniwersalnym pretestem, od którego zależy wybór ANOVA;
jednorodność wariancji też nie wynika z dużej wartości p.

\Needspace{20\baselineskip}
\begin{lstlisting}[language=Julia,caption={Porównanie grup metodą Welcha w Julii: obliczenie statystyki i p-value.}]
using Statistics, Distributions

A = [18, 19, 20, 19, 21, 20, 18, 22]
B = [22, 23, 21, 24, 22, 25, 23, 24]
n, m = length(A), length(B)
vA, vB = var(A), var(B)  # wariancja z mianownikiem n-1
se2 = vA / n + vB / m
tstat = (mean(A) - mean(B)) / sqrt(se2)
df = se2^2 / ((vA/n)^2/(n-1) + (vB/m)^2/(m-1))
pvalue = 2 * cdf(TDist(df), -abs(tstat))
println((t=tstat, df=df, p=pvalue))
\end{lstlisting}

Pakiet \texttt{Distributions.jl} jest zależnością zewnętrzną;
przed uruchomieniem należy go zainstalować w używanym projekcie.
Wynik tego programu jest aproksymacją Satterthwaite'a,
a nie dokładnym rozkładem t przy nierównych wariancjach.

\section{Studium decyzyjne: od wyniku testu do zalecenia}
Rozważmy ponownie trzy technologie z zadania~\ref{zad:anova-three}.
Wynik ANOVA daje bardzo mocne świadectwo przeciw hipotezie
jednakowych średnich, ale sam nie określa najlepszej decyzji.
Przed wyborem technologii należy odpowiedzieć na następujące
pytania.

\begin{enumerate}
\item \emph{Co jest parametrem celu?} Jeśli mierzony czas
pracy chcemy minimalizować, sama statystyczna istotność
różnicy nie przesądza, czy oszczędność czasu rekompensuje
koszty urządzeń. Warto podać różnice średnich w jednostkach
fizycznych, a nie jedynie statystykę $F$.
\item \emph{Jak powstały grupy?} Jeśli te same urządzenia
badano wielokrotnie, nie można ich traktować jako trzech
prób niezależnych. Jeśli technologię przydzielano zależnie
od typu maszyny, nie można przypisać stwierdzonej różnicy
wyłącznie wpływowi technologii.
\item \emph{Jak sprawdzono model?} Wykresy reszt, rozkład
błędów warunkowych i możliwa heterogeniczność wariancji
są ważniejsze niż mechaniczne uruchomienie testu
Shapiro--Wilka osobno dla każdej grupy. Przy nierównych
wariancjach trzeba rozważyć ANOVA Welcha.
\item \emph{Które porównania są celem?} Jeśli wcześniej
zaplanowano porównanie A z B, można analizować ten kontrast.
Jeżeli po obejrzeniu danych wybieramy wszystkie trzy pary,
trzeba kontrolować wielokrotność (np. procedurą Tukeya
w odpowiednim modelu lub Gamesa--Howella przy nierównych
wariancjach).
\item \emph{Jak raportować?} Należy podać liczebności,
średnie i rozproszenie w grupach, sumy kwadratów, stopnie
swobody, wartość $F$ i p-value, miarę efektu oraz
przedziały ufności dla istotnych praktycznie kontrastów.
Dodatkowo trzeba wskazać ograniczenia modelu i możliwe
czynniki zakłócające.
\end{enumerate}

\begin{uwaga}[Brak istotności a równoważność]
Nawet jeśli test różnicy średnich nie odrzuca hipotezy
zerowej, nie wolno mówić, że dwie metody są równoważne.
Aby wykazać równoważność, należy z góry określić
akceptowalny margines $\Delta>0$ oraz przetestować
odpowiednio hipotezę braku równoważności. Omówimy
formalną procedurę dwóch jednostronnych testów (TOST)
w części poświęconej rozszerzeniom testowania hipotez.
\end{uwaga}



% ================================================================
% ROZDZIAL 11 -- PRZEDZIALY UFNOSCI
% ================================================================
